% §6 Exercises
%
% Cross-reference resolution status:
%
% DEFINED here:
%   \label{ex:ch6-1} through \label{ex:ch6-11}
%
% RESOLVED (used here, defined elsewhere):
%   \Cref{prop:union-fpr-fnr}            -- §6.1 (Plan 6 Task 2)
%   \Cref{eq:union-fpr}                  -- §6.1 (Plan 6 Task 2)
%   \Cref{eq:union-fnr}                  -- §6.1 (Plan 6 Task 2)
%   \Cref{prop:intersection-fpr-fnr}     -- §6.2 (Plan 6 Task 3)
%   \Cref{prop:complement-swap}          -- §6.2 (Plan 6 Task 3)
%   \Cref{prop:difference-fpr-fnr}       -- §6.2 (Plan 6 Task 3)
%   \Cref{prop:fnr-zero-closure}         -- §6.4 (Plan 6 Task 5)
%   \Cref{def:perfect-filter-line}       -- §6.4 (Plan 6 Task 5)
%   \Cref{def:one-sided}                 -- chapter 5 §5.4 (Plan 5)
%   \Cref{thm:composition}               -- chapter 3 §3.3 (Plan 3)
%   \Cref{eq:composition-theorem}        -- chapter 3 §3.3 (Plan 3)
%
% Difficulty-marker macros (not in alex.sty; defined locally if not yet defined):
\providecommand{\exstar}{$\star$}
\providecommand{\exstarstar}{$\star\star$}
\providecommand{\exstarstarstar}{$\star\star\star$}

\section*{Exercises}
\addcontentsline{toc}{section}{Exercises}

\begin{enumerate}

% -----------------------------------------------------------------------
% WARM-UP  (★ × 3): exercises 1, 2, 3
% -----------------------------------------------------------------------

\item \exstar
\label{ex:ch6-1}
\textbf{Union rates for two FNR-zero Bloom filters.}
Two FNR-zero Bloom filters on the same universe with the same parameters
have operand rates $\fprate_1 = 0.05, \fnrate_1 = 0$ and $\fprate_2 =
0.10, \fnrate_2 = 0$.
\begin{enumerate}
  \item[(a)] Apply \Cref{eq:union-fpr} of \Cref{prop:union-fpr-fnr} to
    compute the union false-positive rate $\fprate_\cup$. Report to
    four decimal places.
  \item[(b)] Apply \Cref{eq:union-fnr} to compute the union
    false-negative rate $\fnrate_\cup$, noting that every term in the
    weighted sum carries at least one operand $\fnrate$ as a factor.
  \item[(c)] Confirm that the inequality
    $\max(\fprate_1, \fprate_2) \le \fprate_\cup \le \fprate_1 +
    \fprate_2$ holds for the numbers you computed, and name the
    second-order term that accounts for the gap between
    $\fprate_\cup$ and $\fprate_1 + \fprate_2$.
\end{enumerate}

\item \exstar
\label{ex:ch6-2}
\textbf{Intersection rates for the same two filters.}
Take the same operand rates as \Cref{ex:ch6-1}: $\fprate_1 = 0.05,
\fnrate_1 = 0, \fprate_2 = 0.10, \fnrate_2 = 0$.
\begin{enumerate}
  \item[(a)] Apply \Cref{prop:intersection-fpr-fnr} to compute the
    intersection false-positive rate $\fprate_\cap$. Report to four
    decimal places.
  \item[(b)] Compute the intersection false-negative rate
    $\fnrate_\cap$. State why $\fnrate_\cap = 0$ when both operands
    are FNR-zero.
  \item[(c)] Compare the magnitudes of $\fprate_\cap$ and
    $\fprate_\cup$ from \Cref{ex:ch6-1}. By what factor does
    intersection reduce the false-positive rate relative to the
    smaller operand $\fprate_1$, and by what factor does union
    inflate it?
\end{enumerate}

\item \exstar
\label{ex:ch6-3}
\textbf{Complement-swap verification.}
An approximate set $\ASet{A}$ has per-block rates $\fprate = 0.03$ and
$\fnrate = 0.07$.
\begin{enumerate}
  \item[(a)] Apply \Cref{prop:complement-swap} to write the rates
    $(\fprate_\complement, \fnrate_\complement)$ of the complement
    $\complement \ASet{A}$.
  \item[(b)] Apply complement a second time. Verify that the
    twice-complemented rates return to $(0.03, 0.07)$, confirming that
    the complement operation is an involution on the rate plane.
  \item[(c)] In one sentence, name the per-bit precedent for this
    swap that chapter 3 §3.2 already established at the noisy-bit
    level.
\end{enumerate}

% -----------------------------------------------------------------------
% MEDIUM  (★★ × 5): exercises 4, 5, 6, 7, 8
% -----------------------------------------------------------------------

\item \exstarstar
\label{ex:ch6-4}
\textbf{Notebook-tied set-op verification.}
Open the chapter 6 companion notebook
\texttt{python/notebooks/ch06-set-composition.ipynb}. Build two Bloom
filters with $m = 5000$ bits, $n = 500$ inserted random strings each,
and $k = 7$ hash functions, drawn from disjoint inserted sets $S_1$
and $S_2$ so that the operand rates approximate $\fprate_1 \approx
\fprate_2 \approx 0.01$ with $\fnrate_1 = \fnrate_2 = 0$.
\begin{enumerate}
  \item[(a)] Construct the union by bit-wise OR of the two underlying
    bit arrays. Query a pool of $10000$ random strings known to be
    absent from both inserted sets, and record the empirical
    false-positive rate $\hat{\fprate}_\cup$.
  \item[(b)] Compute the closed-form prediction
    $1 - (1 - \hat{\fprate}_1)(1 - \hat{\fprate}_2)$ from
    \Cref{eq:union-fpr}, using empirical operand rates from the
    notebook in place of the targets, and compare to
    $\hat{\fprate}_\cup$. State in one sentence how close the
    empirical and closed-form values agree.
  \item[(c)] Repeat the construction with the operand sets $S_1$ and
    $S_2$ chosen to overlap by $250$ elements. Does the empirical
    union FPR still match the formula, and if so, why does the formula
    not depend on the cardinality of $S_1 \cap S_2$? Refer back to
    the FPR derivation in \Cref{prop:union-fpr-fnr}.
\end{enumerate}

\item \exstarstar
\label{ex:ch6-5}
\textbf{Difference algebra.}
\Cref{prop:difference-fpr-fnr} gave the difference rate formulas
$\fprate_\setminus = \fprate_A \cdot \fnrate_B$ and
$\fnrate_\setminus = 1 - (1 - \fnrate_A)(1 - \fprate_B)$ on the bulk
region. Specialize these to two boundary cases.
\begin{enumerate}
  \item[(a)] Take $\ASet{B}$ to be a perfect filter on the FPR-zero
    axis, so $\fprate_B = 0$ and $\fnrate_B \ge 0$. Substitute into
    the difference FPR formula and confirm that
    $\fprate_\setminus = \fprate_A \cdot \fnrate_B$ is the only
    surviving term. Substitute into the difference FNR formula and
    show that $\fnrate_\setminus = \fnrate_A$, that is, $\ASet{A}
    \setminus \ASet{B}$ inherits its FNR from the first operand alone.
  \item[(b)] Now swap: take $\ASet{B}$ to be FNR-zero, so
    $\fnrate_B = 0$ and $\fprate_B \ge 0$. Substitute into the
    difference FPR formula and confirm $\fprate_\setminus = 0$. Then
    substitute into the difference FNR formula and derive
    $\fnrate_\setminus = 1 - (1 - \fnrate_A)(1 - \fprate_B)$. In one
    sentence, explain the mechanism: false positives of $\ASet{B}$
    become wrongful exclusions in $\ASet{A} \setminus \ASet{B}$, and
    therefore false negatives of the difference.
  \item[(c)] Combine the two boundary readings to argue that the
    difference operation is FPR-preserving when the second operand
    is FNR-zero, and FNR-preserving when the second operand is a
    perfect filter on the FPR-zero axis.
\end{enumerate}

\item \exstarstar
\label{ex:ch6-6}
\textbf{De Morgan duality under noise.}
The classical De Morgan laws assert
$\complement(A \cup B) = \complement A \cap \complement B$ and
$\complement(A \cap B) = \complement A \cup \complement B$. This
exercise verifies that the rate algebra of chapter 6 preserves the
first of these identities under the cross-set independence assumption
of \Cref{prop:union-fpr-fnr}.
\begin{enumerate}
  \item[(a)] Compute the rates of $\complement(\ASet{A} \cup
    \ASet{B})$ by first applying \Cref{prop:union-fpr-fnr} to get the
    union rates and then applying \Cref{prop:complement-swap}. Write
    the resulting $(\fprate, \fnrate)$ as a pair of expressions in the
    operand rates.
  \item[(b)] Compute the rates of $\complement \ASet{A} \cap
    \complement \ASet{B}$ by first applying \Cref{prop:complement-swap}
    to each operand and then applying \Cref{prop:intersection-fpr-fnr}
    to the complements.
  \item[(c)] Compare the two answers and confirm that the FPR
    expressions agree and the FNR expressions agree. The De Morgan
    duality holds for the rate algebra whenever the cross-set
    independence assumption is in force for both compositions.
\end{enumerate}

\item \exstarstar
\label{ex:ch6-7}
\textbf{Independence sensitivity.}
The FPR derivation in the proof of \Cref{prop:union-fpr-fnr} assumed
that the two operand error events at a query point $x$ are mutually
independent. This exercise constructs a small example where the
assumption fails and shows what happens to the formula.
\begin{enumerate}
  \item[(a)] Let $\ASet{A}$ and $\ASet{B}$ be two Bloom filters that
    share \emph{exactly the same} hash family $h_1, \ldots, h_k$,
    that store the same set $S$, and that are derived from the same
    bit array. Argue that the two operand predicates are not merely
    correlated but identical: $\ASet{A}(x) = \ASet{B}(x)$ for every
    $x \in \Universe$.
  \item[(b)] Compute the actual union FPR for this construction.
    Since the two predicates are identical, the OR of their outputs
    agrees with the single-filter output, and the union FPR equals
    the common operand FPR $\fprate$. Compare to the formula
    prediction $1 - (1 - \fprate)(1 - \fprate) = 2 \fprate - \fprate^2$
    and describe the direction of the discrepancy.
  \item[(c)] In one sentence, identify which structural property of
    the construction the formula presupposes, and which paragraph of
    \Cref{sec:union-of-sets} flags the requirement that the bit-wise
    OR construction needs distinct (or at least uncorrelated) hash
    families to match the rate prediction.
\end{enumerate}

\item \exstarstar
\label{ex:ch6-8}
\textbf{Intersection escape from the line.}
\Cref{prop:fnr-zero-closure} stated that intersection of two FNR-zero
sets stays FNR-zero, but the same need not hold when one operand is
FNR-zero and the other is FPR-zero.
\begin{enumerate}
  \item[(a)] Let $\ASet{A}$ be FNR-zero with $\fprate_A > 0$ and
    $\fnrate_A = 0$, and let $\ASet{B}$ be FPR-zero with
    $\fprate_B = 0$ and $\fnrate_B > 0$. Substitute these operand
    rates into \Cref{prop:intersection-fpr-fnr} and derive the
    intersection rates as expressions in $\fprate_A$ and $\fnrate_B$.
  \item[(b)] Verify that $\fprate_\cap = 0$ in this regime, and that
    $\fnrate_\cap = \fnrate_B > 0$. The intersection has migrated to
    the FPR-zero axis: starting from one operand on each axis, the
    composition has landed on a single axis but on the FPR-zero one,
    not the FNR-zero one of $\ASet{A}$.
  \item[(c)] Now perturb the construction: let $\ASet{B}$ instead be
    two-sided with $\fprate_B > 0$ and $\fnrate_B > 0$, leaving
    $\ASet{A}$ FNR-zero. Substitute again and show that both rates
    of the intersection become strictly positive. In one sentence,
    name the geometric phenomenon (the composed point now lies in the
    interior of the rate square) and refer the reader to
    \Cref{def:perfect-filter-line} for the line whose escape this
    exercise concretely exhibits.
\end{enumerate}

% -----------------------------------------------------------------------
% CHALLENGING / EXTENSION  (★★★ × 3): exercises 9, 10, 11
% -----------------------------------------------------------------------

\item \exstarstarstar
\label{ex:ch6-9}
\textbf{Three-way union via the composition theorem.}
Extend \Cref{prop:union-fpr-fnr} to three operands. Let $\ASet{A},
\ASet{B}, \ASet{C}$ be three approximate sets with per-operand FPRs
$\fprate_A, \fprate_B, \fprate_C$, all pairwise independent at every
query point.
\begin{enumerate}
  \item[(a)] Compute $\fprate_{\cup}$ for $\ASet{A} \cup \ASet{B}
    \cup \ASet{C}$ by chained pairwise union: first apply
    \Cref{eq:union-fpr} to $\ASet{A} \cup \ASet{B}$, then apply it
    again to the result composed with $\ASet{C}$. Simplify the
    nested expression to a single product form.
  \item[(b)] Compute $\fprate_{\cup}$ directly from independence as
    in the proof of \Cref{prop:union-fpr-fnr}: a false positive of
    the three-way union at an $x$ outside all three sets requires at
    least one of the three operand predicates to report $\top$, so
    the union FPR is $1$ minus the joint probability that all three
    report $\bot$. Use independence to factor the joint into a
    product.
  \item[(c)] Confirm that (a) and (b) agree. Then argue that the
    result is the composition theorem \Cref{thm:composition} of
    chapter 3 §3.3 instantiated at $k = 3$, that is,
    \Cref{eq:composition-theorem} with $\eta_i = \fprate_i$. The
    same product formula governs every $k \ge 2$ by induction.
\end{enumerate}

\item \exstarstarstar
\label{ex:ch6-10}
\textbf{Perfect filter as algebraic inverse.}
A perfect filter on a designated input set $A$ is an approximate set
with $\fprate = 0$ on $A$, paying with a positive $\fnrate$ on the
complement domain. \Cref{prop:fnr-zero-closure} part (4) suggested
that pairing such a filter with an FNR-zero filter can give an
exact-set test on a constrained domain. This exercise sketches the
construction.
\begin{enumerate}
  \item[(a)] Let $\ASet{A}_+$ be a perfect filter for $A$ with rates
    $(0, \fnrate_+)$ on its designed domain, and let $\ASet{A}_-$ be
    an FNR-zero filter for $A$ with rates $(\fprate_-, 0)$. Construct
    the composed predicate $\ASet{A}_+(x) \land \ASet{A}_-(x)$ and
    apply \Cref{prop:intersection-fpr-fnr} to compute its rates.
    Confirm $\fprate = \fprate_- \cdot 0 = 0$ and
    $\fnrate = 1 - (1 - \fnrate_+)(1 - 0) = \fnrate_+$.
  \item[(b)] State the condition under which $\fnrate_+ = 0$ also.
    The perfect filter has zero FNR on a subset of its designed
    domain (its ``acceptance set''); if that subset coincides with
    $A$ itself, the composition has both rates zero on $A$, that is,
    an exact-set test.
  \item[(c)] Describe the architectural reading: the FNR-zero filter
    serves as a fast first-pass classifier (eliminating most
    non-members with one cheap query), and the perfect filter
    serves as a precise second-pass classifier (verifying the
    remaining candidates with certainty on the designed domain).
    Sketch why this two-pass composition is the algorithmic skeleton
    of the perfect hash filter that chapter 15 will develop in
    detail.
\end{enumerate}

\item \exstarstarstar
\label{ex:ch6-11}
\textbf{Entropy floor on the union operation.}
Chapter 8 will prove that any FNR-zero approximate set at false-positive
rate $\fprate$ requires at least $-\log_2 \fprate$ bits per stored
element to represent, and chapter 14 will exhibit the Bernoulli Hash
Function construction that achieves this floor. This exercise applies
the floor to the union of two FNR-zero sets and shows that no clever
implementation of union can beat it.
\begin{enumerate}
  \item[(a)] Let $\ASet{A}$ and $\ASet{B}$ be two FNR-zero approximate
    sets for $A, B \subseteq \Universe$, with operand FPRs
    $\fprate_A, \fprate_B$. By \Cref{prop:fnr-zero-closure} part (1)
    the union $\ASet{A} \cup \ASet{B}$ is FNR-zero with FPR
    $\fprate_{\cup} = 1 - (1 - \fprate_A)(1 - \fprate_B)$. Apply the
    chapter 8 floor (anticipated) to the union: any representation
    of $\ASet{A} \cup \ASet{B}$ as a single FNR-zero approximate set
    of the objective set $A \cup B$ requires at least $|A \cup B|
    \cdot (-\log_2 \fprate_{\cup})$ bits.
  \item[(b)] Argue that the chapter 6 composition formulas do not
    cheat the floor. The composed structure stores no more than the
    two operands together (its bit array, in the Bloom case, is the
    OR of the two operand bit arrays and so occupies $m$ bits, the
    same as each operand). However, $\fprate_{\cup} > \max(\fprate_A,
    \fprate_B)$, so the floor for the composed structure permits a
    correspondingly looser bound; the bit budget allocated by each
    operand at its own floor is sufficient (and asymptotically tight,
    by chapter 14) for the composed FPR.
  \item[(c)] In one or two sentences, draw the structural conclusion
    that motivates chapter 14. The union operation does not save
    bits beyond what the operands already paid for; constructive
    realizations that achieve the floor on each operand also achieve
    the floor on their union, modulo the small inflation $\fprate_{\cup}
    > \max(\fprate_A, \fprate_B)$ that the union FPR formula records.
    Chapter 14 develops the constructive realization (the BHF) that
    saturates this regime.
\end{enumerate}

\end{enumerate}
